\section{Case Study: Self-Evolving Verification in a Programmatic Verifier}
\label{app:self-evolving-verification-case-study}

This appendix provides a concrete example of the self-evolving verification
layer described in Section~\ref{sec:self-evolving-verification}. The goal of this
layer is to use execution-grounded feedback to refine the verification stack and
improve future task synthesis. The example below illustrates a common failure
mode: the agent completed the task, but the verifier queried an outdated
application schema and therefore incorrectly marked several satisfied criteria as
failed. By comparing the programmatic verdict against an LLM reference judgment
on a fixed trajectory, the system identifies the verifier-side defect and updates
the verification logic accordingly.

\paragraph{Task.}
We use the darktable task \texttt{darktable\_batch\_rate\_and\_tag}. The agent is
instructed to import three images, create the tag \texttt{batch\_processed},
attach the tag to all three images, and assign ratings of one, three, and five
stars to \texttt{img\_001.png}, \texttt{img\_002.png}, and
\texttt{img\_003.png}, respectively. The recorded run was produced by
\texttt{kimi-k2.6} and completed in 53 interaction steps. The trajectory was then
frozen and reused throughout the self-evolution procedure.

\paragraph{Reference judgment.}
An LLM judge inspected the full trajectory, post-action screenshots, and final
state. It judged all ten criteria as satisfied: the three images were imported,
the tag \texttt{batch\_processed} was visible and attached to all images, and the
final star ratings matched the requested values. In particular, the judge found
that the tag was created while all three images were selected, and that the final
state showed the expected image information and rating flags for each image.

\begin{table}[H]
\centering
\small
\begin{tabular}{llccc}
\toprule
Round & Source & Passed & Failed & Divergences \\
\midrule
0 & Programmatic verifier before evolution & 6 & 4 & 4 \\
0 & LLM reference judgment & 10 & 0 & -- \\
1 & Programmatic verifier after evolution & 10 & 0 & 0 \\
\bottomrule
\end{tabular}
\caption{
Self-evolution outcome for \texttt{darktable\_batch\_rate\_and\_tag}. The initial
verifier incorrectly failed four tag-related criteria. After updating the
verification logic using execution-grounded feedback, the programmatic verdict
agreed with the LLM reference on all ten criteria.
}
\label{tab:darktable-self-evolving-outcome}
\end{table}

\paragraph{Detected disagreement.}
The comparator found four disagreements between the LLM reference judgment and
the programmatic verifier. All four involved tag state: whether the tag
\texttt{batch\_processed} existed, and whether it was attached to each of the
three imported images. In each case, the LLM judge returned \textsc{True}, while
the verifier returned \textsc{False}. The comparator classified all four
disagreements as verifier-side failures, meaning that the task had been completed
but the checker misjudged the final state.

\begin{table}[H]
\centering
\small
\caption{
Criterion-level disagreements before self-evolution. All failures share the
same root cause: the verifier checked tag metadata using an outdated database
assumption.
}
\begin{tabular}{clccc}
\toprule
Criterion & Description & Verifier & Judge & Classification \\
\midrule
3 & Tag \texttt{batch\_processed} exists & False & True & \texttt{verifier\_wrong} \\
4 & \texttt{img\_001} has tag & False & True & \texttt{verifier\_wrong} \\
5 & \texttt{img\_002} has tag & False & True & \texttt{verifier\_wrong} \\
6 & \texttt{img\_003} has tag & False & True & \texttt{verifier\_wrong} \\
\bottomrule
\end{tabular}
\label{tab:darktable-self-evolving-disagreements}
\end{table}

\paragraph{Root cause.}
All four failures were caused by a single verifier bug. The darktable verifier
assumed that the table \texttt{tags} lived in \texttt{library.db}. In the current
darktable state, however, tag definitions are stored in \texttt{data.db}, while
image-tag associations remain in \texttt{library.db}. As a result, tag-related
SQL queries failed with a missing-table error and were counted as negative
verifier results, even though the final application state contained the expected
tag assignments.

\paragraph{Verifier evolution.}
The self-evolving layer was allowed to modify only the verifier implementation
and documentation, not the agent trajectory, sandbox state, task specification,
or expected output. The verifier update made three changes. First,
\texttt{check\_tag\_exists} and the corresponding tag-query endpoint were
rerouted to query \texttt{data.db}. Second, the image-tag checker was rewritten
to join \texttt{library.db}'s \texttt{tagged\_images} table with
\texttt{data.db}'s \texttt{tags} table. Third, the verifier documentation was
updated to reflect the actual darktable schema. These changes preserve the same
public checker interface while aligning the internal inspection logic with the
real application state.

\begin{table}[H]
\centering
\small
\caption{
Summary of the verifier evolution. The public checker interface was unchanged;
only the internal SQL source and join path were updated.
}
\begin{tabular}{p{0.28\linewidth}p{0.31\linewidth}p{0.31\linewidth}}
\toprule
Check & Before evolution & After evolution \\
\midrule
Tag existence &
Query \texttt{main.tags} inside \texttt{library.db}. &
Query \texttt{main.tags} inside \texttt{data.db}. \\
\addlinespace
Image-tag assignment &
Join \texttt{main.tagged\_images} with \texttt{main.tags} inside \texttt{library.db}. &
Join \texttt{main.tagged\_images} from \texttt{library.db} with \texttt{data.tags} from attached \texttt{data.db}. \\
\bottomrule
\end{tabular}
\label{tab:darktable-self-evolving-update}
\end{table}

\paragraph{Before and after.}
The tag-existence check required only changing the database queried by the
existing SQL statement:
\begin{verbatim}
# Before
rows = _query_sqlite(LIBRARY_DB, sql, (tag_name, f"%|{tag_name}"))

# After
rows = _query_sqlite(DATA_DB, sql, (tag_name, f"%|{tag_name}"))
\end{verbatim}

The image-tag checker required a cross-database join:
\begin{verbatim}
# Before
SELECT t.id AS tag_id, t.name AS tag_name
FROM main.tagged_images ti
JOIN main.tags t ON ti.tagid = t.id
WHERE ti.imgid = ? AND (t.name = ? OR t.name LIKE ?)

# After
SELECT t.id AS tag_id, t.name AS tag_name
FROM main.tagged_images ti
JOIN data.tags t ON ti.tagid = t.id
WHERE ti.imgid = ? AND (t.name = ? OR t.name LIKE ?)
\end{verbatim}

\paragraph{Outcome.}
After self-evolution, the verifier was re-executed on the same cached final
state. The updated verifier passed all ten criteria and had zero remaining
divergences from the LLM reference judgment. This example demonstrates how the
self-evolving verification layer provides an additional feedback channel for the
synthesis pipeline: it identifies brittle verifier assumptions, such as
application schema drift, updates the executable inspection logic, and records
which application states require more careful grounding in future task
generation. In this way, OpenComputer improves its verifier stack over time
while preserving the core principle that agent performance is scored by
executable, application-grounded checks.